
This work investigates whether the coefficient of variation (CV) of linear probe accuracy trajectories can distinguish spurious from core features in pretrained representations. The hypothesis posits that spurious features emerge uniformly across sample subsets (low CV) while core features emerge differentially (high CV). Experiments on the Waterbirds dataset using frozen CLIP ViT-B/16 features yield AUC = 0.0 for CV-based classification of spurious versus core features. Both background (spurious) and bird type (core) features exhibit nearly identical CV values (0.0393 and 0.0360, respectively), providing no discriminative signal. The failure stems from feature saturation: the pretrained model has already learned both concepts, eliminating emergence dynamics. Linear probes converge immediately at all regularization strengths, producing flat trajectories. This negative result clarifies that emergence uniformity is a training-time phenomenon that cannot be recovered from frozen pretrained features and motivates future work toward training-time measurement with unfrozen representations.

